\chapter{Initial Feasibility Study}

\section{Introduction}

This chapter presents an initial feasibility study for the major problems identified in Chapter 1 of the Information System Analysis report on NextStepBD. The purpose of this chapter is to evaluate whether the proposed improvements are technically possible, economically reasonable, and operationally acceptable for the organization.

The feasibility analysis follows the standard Technical, Economic, and Operational (TEO) framework. Each problem is examined from three perspectives:

\begin{itemize}
    \item \textbf{Technical Feasibility:} Whether the proposed solution can be implemented using available technology, tools, and technical skills.
    \item \textbf{Economic Feasibility:} Whether the expected benefits justify the cost of implementation, training, and maintenance.
    \item \textbf{Operational Feasibility:} Whether the organization can realistically adopt the solution within its existing workflow, culture, and manpower.
\end{itemize}

NextStepBD is a software development, automation, and consultancy-based company. The company currently uses tools such as WhatsApp, Slack, Email, GitHub, Microsoft Excel, and Google Sheets for communication, development, reporting, and coordination. Although these tools support daily operations, the absence of a centralized and integrated information system creates inefficiency, duplication, tracking difficulty, and dependency on manual coordination.

The goal of this feasibility study is to identify which improvement initiatives should be prioritized so that NextStepBD can build a more structured, scalable, and efficient operational system.

\section{Initial Feasibility Analysis}

\subsection{Problem 01: Fragmented System and Lack of Integration}

\subsubsection{Statement of the Problem}

NextStepBD currently uses multiple disconnected platforms such as WhatsApp, GitHub, Google Sheets, Excel, Slack, and Email. Each tool performs a useful function, but there is no centralized system that connects communication, task management, development progress, client requirements, and reporting.

As a result, information is scattered across different platforms. Employees may need to manually copy data from one system to another, which increases the chance of mistakes, duplication, and delays. Managers also face difficulty in obtaining a single, accurate view of project status.

\subsubsection{Summary of Findings}

Introducing an integrated workflow management system is technically feasible because many cloud-based platforms already support integration with GitHub, Google Sheets, Slack, Email, and other business tools. The solution is economically feasible because it can reduce manual effort, rework, and coordination delays. Operationally, the solution is feasible if implemented gradually, starting with one or two pilot projects before expanding across the whole organization.

\subsubsection{Details of Findings}

\begin{itemize}
    \item \textbf{Technical Feasibility:} Tools such as Jira, ClickUp, Notion, Trello, Zoho Projects, and Monday.com can integrate task tracking, document storage, reporting, and communication. These platforms can also connect with GitHub and Google Workspace. Therefore, the company does not need to develop the entire system from the beginning.
    
    \item \textbf{Economic Feasibility:} The estimated annual cost of a cloud-based integrated system for a small to medium-sized team may range from BDT 3--8 lakhs, depending on the number of users and selected platform. The cost is justified because the system can reduce manual coordination, reporting time, and project delays.
    
    \item \textbf{Operational Feasibility:} Since NextStepBD already uses digital tools, employees are familiar with online platforms. However, training and clear usage policies will be required. The system should be introduced gradually so that employees do not feel sudden pressure.
    
    \item \textbf{Risk Assessment:} The main risk is resistance from employees who are used to existing informal tools. This can be reduced through training, management support, and a phased implementation plan.
\end{itemize}

\subsubsection{Recommendation}

NextStepBD should adopt a centralized project and workflow management platform within 3--4 months. The system should integrate task assignment, GitHub development updates, client requirements, reporting, and internal communication. A pilot implementation should be conducted with selected projects before full organizational deployment.

\begin{table}[H]
\centering
\caption{Integrated System Implementation -- TEO Feasibility Summary}
\begin{tabular}{|p{3cm}|p{2.5cm}|p{7cm}|}
\hline
\textbf{Dimension} & \textbf{Rating} & \textbf{Justification} \\
\hline
Technical & High & Cloud-based integrated platforms are available and can connect with existing tools such as GitHub, Google Sheets, Slack, and Email. \\
\hline
Economic & High & Reduces manual work, duplication, reporting delay, and coordination cost. \\
\hline
Operational & Medium & Requires training and gradual adoption by employees and managers. \\
\hline
Overall Viability & Feasible & Recommended for early implementation. \\
\hline
\end{tabular}
\end{table}

\subsection{Problem 02: Over-Reliance on Informal Communication Channels}

\subsubsection{Statement of the Problem}

WhatsApp is heavily used by NextStepBD for client communication and internal coordination. While WhatsApp provides fast communication, it is not designed for structured project management, official documentation, task assignment, or decision tracking.

Important messages, client requirements, and project decisions may become buried inside long chat threads. This creates difficulty in searching past discussions, identifying responsible persons, and maintaining accountability.

\subsubsection{Summary of Findings}

A structured communication and documentation policy is feasible for NextStepBD. The company does not need to stop using WhatsApp completely; instead, WhatsApp can be used for quick communication while formal decisions, requirements, task assignments, and approvals should be recorded in a centralized project system.

\subsubsection{Details of Findings}

\begin{itemize}
    \item \textbf{Technical Feasibility:} Slack, Microsoft Teams, Google Chat, and project management tools can organize communication into channels, threads, project spaces, and searchable records. These tools can also integrate with task boards and GitHub.
    
    \item \textbf{Economic Feasibility:} The company can begin with free or low-cost plans. Since Slack is already used for international collaboration, better use of Slack or a similar tool would not require a major additional investment.
    
    \item \textbf{Operational Feasibility:} Employees are already comfortable with chat-based communication. Therefore, shifting important project communication from informal WhatsApp groups to structured project channels is realistic. However, management must define which type of communication belongs to which platform.
    
    \item \textbf{Risk Assessment:} Employees may continue using WhatsApp out of habit. This can be controlled by setting rules such as: urgent messages may be sent on WhatsApp, but requirements, approvals, and task updates must be documented in the official system.
\end{itemize}

\subsubsection{Recommendation}

NextStepBD should introduce a formal communication policy. WhatsApp should remain available for quick notifications and emergency communication, but official client requirements, decisions, approvals, and task updates should be recorded in the project management system or Slack channels.

\begin{table}[H]
\centering
\caption{Communication System -- Current State vs Proposed State}
\begin{tabular}{|p{4cm}|p{5cm}|p{5cm}|}
\hline
\textbf{Area} & \textbf{Current State} & \textbf{Proposed State} \\
\hline
Client Communication & Mostly WhatsApp and Email & Structured client communication log with Email and project portal \\
\hline
Internal Discussion & WhatsApp groups & Slack or project-based channels \\
\hline
Decision Tracking & Difficult to retrieve & Recorded in centralized system \\
\hline
Task Assignment & Informal messages & Assigned with deadline and responsible person \\
\hline
Documentation & Scattered & Searchable and centralized \\
\hline
\end{tabular}
\end{table}

\subsection{Problem 03: Unstructured Workflow and Weak Requirement Management}

\subsubsection{Statement of the Problem}

Client requirements at NextStepBD are often received through WhatsApp messages, shared documents, or informal discussions. Because there is no standardized requirement collection format, requirements may be incomplete, unclear, or frequently changed.

This causes confusion among developers and managers. Developers may need to assume missing details, which increases the risk of rework, project delay, and mismatch between client expectations and final delivery.

\subsubsection{Summary of Findings}

A structured requirement management process is highly feasible. The company can introduce standard requirement templates, client approval checkpoints, and change request forms. These improvements do not require heavy infrastructure but can significantly improve project clarity and delivery quality.

\subsubsection{Details of Findings}

\begin{itemize}
    \item \textbf{Technical Feasibility:} Requirement management can be implemented using Google Forms, Notion, Jira, Trello, ClickUp, or a custom internal portal. These tools can store requirement descriptions, attachments, priority levels, acceptance criteria, and approval status.
    
    \item \textbf{Economic Feasibility:} The cost is low because the company already uses Google Workspace and other digital tools. The main investment will be time spent designing templates and training employees.
    
    \item \textbf{Operational Feasibility:} Managers and developers will benefit directly from clearer requirements. Clients may also appreciate a structured process because it reduces misunderstanding. However, the process should be simple so clients do not feel overloaded.
    
    \item \textbf{Risk Assessment:} If the requirement form is too complicated, clients and managers may avoid using it. Therefore, the format should be concise and practical.
\end{itemize}

\subsubsection{Recommendation}

NextStepBD should create a standard Software Requirement Specification (SRS) template or lightweight requirement form. Every new project or feature request should include objective, user role, functional requirement, non-functional requirement, priority, deadline, and approval status. Requirement changes should be handled through a formal change request process.

\begin{table}[H]
\centering
\caption{Proposed Requirement Management Template}
\begin{tabular}{|p{4cm}|p{9cm}|}
\hline
\textbf{Field} & \textbf{Purpose} \\
\hline
Requirement ID & Unique identification for tracking \\
\hline
Client Name & Identifies the source of the requirement \\
\hline
Requirement Description & Explains what needs to be developed \\
\hline
Priority & High, Medium, or Low \\
\hline
Acceptance Criteria & Defines when the requirement will be considered complete \\
\hline
Responsible Person & Assigns ownership \\
\hline
Approval Status & Pending, Approved, Revised, or Rejected \\
\hline
Change History & Records updates and modifications \\
\hline
\end{tabular}
\end{table}

\subsection{Problem 04: Absence of Centralized Task and Project Management System}

\subsubsection{Statement of the Problem}

Task assignment and project progress tracking are currently handled mostly through manual communication. Managers assign tasks through WhatsApp or direct messages, while developers use GitHub mainly for code version control. There is no single dashboard for tracking deadlines, dependencies, progress, bottlenecks, and responsibility.

This makes it difficult for managers to monitor all tasks efficiently. Developers may also lack clarity about priority, deadline, and dependency.

\subsubsection{Summary of Findings}

A centralized task and project management system is technically and economically feasible. Since the company already uses GitHub, a project management system can be integrated with GitHub issues, pull requests, and development workflow. This will improve visibility, accountability, and delivery planning.

\subsubsection{Details of Findings}

\begin{itemize}
    \item \textbf{Technical Feasibility:} GitHub Projects, Jira, ClickUp, Trello, and Zoho Projects can provide task boards, Kanban views, sprint planning, deadlines, and assignment tracking. GitHub Projects may be especially suitable because the development team already uses GitHub.
    
    \item \textbf{Economic Feasibility:} GitHub Projects can be used with minimal additional cost. More advanced platforms may require subscription fees, but the benefit of improved project control can justify the cost.
    
    \item \textbf{Operational Feasibility:} Developers are already familiar with GitHub, so adopting GitHub Issues or GitHub Projects would be easier than introducing a completely new system. Managers will need training on dashboard monitoring and task reporting.
    
    \item \textbf{Risk Assessment:} If tasks are not updated regularly, the dashboard will become unreliable. Therefore, task updates should be part of the daily workflow.
\end{itemize}

\subsubsection{Recommendation}

NextStepBD should implement a centralized task management system using GitHub Projects or a similar project management tool. Each task should have a responsible person, deadline, priority, status, and related GitHub issue or pull request. Managers should review the dashboard regularly to identify delays and bottlenecks.

\begin{table}[H]
\centering
\caption{Task Management System -- TEO Feasibility Summary}
\begin{tabular}{|p{3cm}|p{2.5cm}|p{7cm}|}
\hline
\textbf{Dimension} & \textbf{Rating} & \textbf{Justification} \\
\hline
Technical & High & GitHub Projects and other project tools are available and suitable for software teams. \\
\hline
Economic & High & Low-cost implementation; improves productivity and reduces delay. \\
\hline
Operational & High & Developers already use GitHub, so adoption is practical. \\
\hline
Overall Viability & Highly Feasible & Recommended for immediate implementation. \\
\hline
\end{tabular}
\end{table}

\subsection{Problem 05: Inefficient Data Management and Reporting}

\subsubsection{Statement of the Problem}

NextStepBD relies heavily on Microsoft Excel and Google Sheets for reporting, financial data handling, client-side collaboration, and operational records. Although these tools are useful for basic data handling, they become inefficient when data volume increases or when multiple people update files manually.

Data may become scattered across different sheets, resulting in inconsistency, duplication, and difficulty in generating real-time reports.

\subsubsection{Summary of Findings}

Improving the reporting system is feasible through a centralized database and dashboard-based reporting. The company can start with a simple system that collects project, client, sales, and operational data in one place and generates automated reports.

\subsubsection{Details of Findings}

\begin{itemize}
    \item \textbf{Technical Feasibility:} Google Looker Studio, Microsoft Power BI, Airtable, Zoho Analytics, or a custom web-based dashboard can be used. These tools can connect with spreadsheets, databases, and project management systems.
    
    \item \textbf{Economic Feasibility:} The initial cost can be controlled by using existing Google Sheets with dashboard tools. A fully custom reporting system may cost more, but it can be developed gradually by the internal team.
    
    \item \textbf{Operational Feasibility:} Managers and administrative staff already use spreadsheets, so dashboard-based reporting would be easier to adopt. The main change will be maintaining clean and standardized data.
    
    \item \textbf{Risk Assessment:} Poor data entry discipline can reduce report accuracy. The company should define standard data formats, access control, and validation rules.
\end{itemize}

\subsubsection{Recommendation}

NextStepBD should introduce a centralized reporting dashboard. At the first stage, Google Sheets can remain as the data source, but reports should be automated using dashboard tools. In the long term, the company should move important operational data to a centralized database.

\begin{table}[H]
\centering
\caption{Reporting Improvement -- Estimated Cost-Benefit Analysis}
\begin{tabular}{|p{5cm}|p{4cm}|p{4cm}|}
\hline
\textbf{Item} & \textbf{Estimated Annual Cost} & \textbf{Expected Benefit} \\
\hline
Dashboard Tool Setup & BDT 1--2 Lakhs & Faster reporting \\
\hline
Data Cleaning and Standardization & BDT 1 Lakh & Improved accuracy \\
\hline
Training & BDT 0.5 Lakh & Better user adoption \\
\hline
Manual Reporting Time Reduction & -- & Saves manager and admin time \\
\hline
Error Reduction & -- & Improves decision-making quality \\
\hline
\end{tabular}
\end{table}

\subsection{Problem 06: Scalability and Dependency Issues}

\subsubsection{Statement of the Problem}

As NextStepBD grows, its current workflow creates dependency on specific individuals, especially managers who act as intermediaries between clients and developers. If these key individuals are unavailable, task assignment, clarification, reporting, and decision-making may be delayed.

The lack of standardized processes also makes it difficult to onboard new employees and manage larger teams.

\subsubsection{Summary of Findings}

Scalability can be improved by standardizing workflow, documenting processes, and introducing role-based access and responsibility. This is technically feasible because process documentation and workflow automation tools are readily available. It is economically justified because it reduces dependency risk and improves continuity.

\subsubsection{Details of Findings}

\begin{itemize}
    \item \textbf{Technical Feasibility:} Standard operating procedures, onboarding checklists, knowledge bases, and workflow automation can be implemented using Notion, Confluence, Google Drive, GitHub Wiki, or a custom internal portal.
    
    \item \textbf{Economic Feasibility:} The cost is moderate because the main investment is process design, documentation, and training. The benefit includes faster onboarding, fewer delays, and reduced dependency on individual managers.
    
    \item \textbf{Operational Feasibility:} The flat organizational structure of NextStepBD supports quick decision-making, but growth requires more formal role definition. Employees may accept the change if it reduces confusion and improves work clarity.
    
    \item \textbf{Risk Assessment:} Documentation may become outdated if not maintained. A responsible person or team should review process documents regularly.
\end{itemize}

\subsubsection{Recommendation}

NextStepBD should create a centralized knowledge base and standard operating procedures for project intake, requirement analysis, task assignment, reporting, client communication, and employee onboarding. The company should also define backup roles so that work does not depend on one person only.

\begin{table}[H]
\centering
\caption{Scalability Improvement Plan}
\begin{tabular}{|p{4cm}|p{5cm}|p{5cm}|}
\hline
\textbf{Area} & \textbf{Current Issue} & \textbf{Proposed Improvement} \\
\hline
Manager Dependency & Managers handle many coordination tasks manually & Use workflow system and assign backup responsibility \\
\hline
Employee Onboarding & Informal knowledge transfer & Create onboarding checklist and knowledge base \\
\hline
Process Documentation & Limited or scattered documentation & Maintain SOPs in a centralized repository \\
\hline
Decision Continuity & Delays when key persons are unavailable & Define role-based approval and escalation process \\
\hline
\end{tabular}
\end{table}

\section{Comparative Feasibility Summary}

The following table summarizes the feasibility of the six identified problems and the recommended priority level for implementation.

\begin{table}[H]
\centering
\caption{Comparative Feasibility Analysis -- All Six Problems}
\begin{tabular}{|c|p{5cm}|c|c|c|p{2.5cm}|}
\hline
\textbf{No.} & \textbf{Problem} & \textbf{Technical} & \textbf{Economic} & \textbf{Operational} & \textbf{Priority} \\
\hline
01 & Fragmented System and Lack of Integration & High & High & Medium & Urgent \\
\hline
02 & Over-Reliance on Informal Communication Channels & High & High & Medium & High \\
\hline
03 & Unstructured Workflow and Weak Requirement Management & High & High & High & Urgent \\
\hline
04 & Absence of Centralized Task and Project Management System & High & High & High & Urgent \\
\hline
05 & Inefficient Data Management and Reporting & High & Medium & Medium & High \\
\hline
06 & Scalability and Dependency Issues & Medium & High & Medium & Medium \\
\hline
\end{tabular}
\end{table}

From the above analysis, the most urgent improvements are the implementation of a centralized task/project management system, structured requirement management, and integration of the currently disconnected tools. These initiatives directly affect daily productivity, project delivery, accountability, and client satisfaction.

\section{Recommended Implementation Roadmap}

The proposed solutions should not be implemented all at once. A phased approach is more realistic and less disruptive for the organization.

\begin{table}[H]
\centering
\caption{Proposed Implementation Roadmap}
\begin{tabular}{|p{3cm}|p{7cm}|p{4cm}|}
\hline
\textbf{Phase} & \textbf{Activities} & \textbf{Estimated Duration} \\
\hline
Phase 1 & Select project management platform, create task board, define requirement template & Month 1--2 \\
\hline
Phase 2 & Pilot implementation with selected projects and train managers/developers & Month 3--4 \\
\hline
Phase 3 & Integrate GitHub, communication channels, and reporting dashboards & Month 5--6 \\
\hline
Phase 4 & Create knowledge base, SOPs, onboarding guide, and backup role structure & Month 7--9 \\
\hline
Phase 5 & Review performance, collect feedback, and optimize system usage & Month 10--12 \\
\hline
\end{tabular}
\end{table}

\section{Conclusion}

The initial feasibility study shows that the problems identified in NextStepBD's current operational system are solvable through practical and cost-effective information system improvements. The company does not need to replace all existing tools immediately. Instead, it should connect and organize them through a centralized workflow, task management, requirement management, communication, and reporting structure.

The highest priority should be given to centralized task management, structured requirement collection, and tool integration because these directly reduce confusion, manual coordination, and project delays. Communication policy and reporting automation should be implemented next, followed by process documentation and scalability improvements.

Overall, the proposed improvements are technically feasible, economically justified, and operationally realistic. If implemented in phases, the improved system will help NextStepBD achieve better coordination, stronger accountability, faster reporting, improved client satisfaction, and sustainable organizational growth.